\section{实验设计}
实验目标是从功能可用性、系统性能、问答质量与风险预测效果四个维度评估平台。实验环境采用Ubuntu服务器，后端服务容器化部署，前端通过内网访问。数据集由脱敏病历样本、随访记录与公开医学知识条目构成，按时间划分训练集、验证集与测试集。

\section{评价指标}
问答模块采用准确率、相关性评分与人工可用性评分联合评价；风险预测模块采用准确率（Accuracy）、召回率（Recall）、F1值与AUC作为主要指标；系统性能采用P95响应时延与并发吞吐量。通过多指标联合评估，避免单一指标造成的偏差。

\section{实验结果与分析}
在当前原型配置下，系统已完成端到端流程验证：数据可稳定入库，问答可返回带引用片段的答案，风险服务可输出分层结果并触发告警。初步测试显示，问答模块在常见咨询场景下具有较高可用性，风险模型对高危样本具备较好的召回能力。后续将进一步补充多中心样本上的严格对比实验，重点完善以下内容：
\begin{itemize}
  \item 与传统规则法、单模型法的横向对比结果；
  \item 不同阈值下精确率-召回率权衡；
  \item 特征消融实验与错误案例分析；
  \item 模型漂移监测与周期性重训练策略验证。
\end{itemize}

\section{小结}
实验表明该平台具备“可部署、可使用、可扩展”的工程基础。下一阶段将通过更大规模真实世界数据验证模型泛化能力与临床实用价值，并完善伦理合规与安全审计流程。